
          %%%%% ~~~~~~~~~~~~~~~~~~~~ %%%%%

\chapter{Classification of irreducible representations of semisimple Lie algebras via highest weights}
\setcounter{theorem}{0}
\setcounter{equation}{0}

%\cite{vanDerVaart1996}
%\cite{Kosorok2008}

%\renewcommand{\theenumi}{\alph{enumi}}
%\renewcommand{\labelenumi}{\textnormal{(\theenumi)}$\;\;$}
\renewcommand{\theenumi}{\roman{enumi}}
\renewcommand{\labelenumi}{\textnormal{(\theenumi)}$\;\;$}

          %%%%% ~~~~~~~~~~~~~~~~~~~~ %%%%%

\section{Semisimplicity $\Longleftrightarrow$ direct sum of simple Lie algebras}

          %%%%% ~~~~~~~~~~~~~~~~~~~~ %%%%%

\vskip 0.5cm
\begin{definition}[Irreducibility and simplicity of Lie algebras]
\mbox{}
\vskip 0.1cm
\noindent
A Lie algebra \,$\mathfrak{g}$\, is said to be
\begin{itemize}
\item
    \textbf{irreducible} if the only ideals in \,$\mathfrak{g}$\,
    are \,$\mathfrak{g}$\, and \,$\{\,0\,\}$\,, and
\item
    \textbf{simple} if \,$\mathfrak{g}$\, is irreducible and
    \,{\color{red}$\dim(\,\mathfrak{g}\,) \,\geq\, 2$}.\,
\end{itemize}
\end{definition}

          %%%%% ~~~~~~~~~~~~~~~~~~~~ %%%%%

\begin{definition}[Reductiveness and semisimplicity of complex Lie algebras]
\mbox{}
\vskip 0.1cm
\noindent
A complex Lie algebra \,$\mathfrak{g}$\, is said to be
\begin{itemize}
\item
    \textbf{reductive} if there exists a real Lie subalgebra
    \,$\mathfrak{k} \,\subset\, \mathfrak{g}$\,
    such that
    \,$\mathfrak{g} \,\cong\, \mathfrak{k}_{\C}$,\,
    where \,$\mathfrak{k}$\, is isomorphic
    to the Lie algebra of a {\color{red}compact} matrix Lie group, and
\item
    \textbf{semisimple} if \,$\mathfrak{g}$\, is reductive and
    the centre of \,$\mathfrak{g}$\, is trivial.
\end{itemize}
\end{definition}

          %%%%% ~~~~~~~~~~~~~~~~~~~~ %%%%%

\begin{definition}[Compact real forms of semisimple complex Lie algebras]
\mbox{}
\vskip 0.1cm
\noindent
Suppose \,$\mathfrak{g}$\, is a semisimple complex Lie algebra.
A real Lie subalgebra \,$\mathfrak{k} \,\subset \mathfrak{g}\,$\,
is called a \textbf{compact real form} of \,$\mathfrak{g}$\, if
\,$\mathfrak{k}$\, is isomorphic to the Lie algebra of some
{\color{red}compact} matrix Lie group, and, for each \,$Z \in \mathfrak{g}$,\,
there exist unique \,$X,\, Y \in \mathfrak{k}$\, such that
\,$Z \,=\, X \,+\, \i\,Y$.\,
\end{definition}

          %%%%% ~~~~~~~~~~~~~~~~~~~~ %%%%%

After first establishing a number of auxiliary results, we will prove the following:

          %%%%% ~~~~~~~~~~~~~~~~~~~~ %%%%%

\vskip 0.5cm
\begin{theorem}[Semisimple $\Longrightarrow$ direct sum of simple Lie algebras]
\label{SemisimpleImpliesDirectSumOfSimples}
\mbox{}
\vskip 0.1cm
\noindent
If \,$\mathfrak{g}$\, is a semisimple complex Lie algebra,
then \,$\mathfrak{g}$\, can be expressed as a Lie algebra direct sum
\begin{equation*}
\mathfrak{g}
\;\; = \;\;
    \overset{m}{\underset{j = 1}{\bigoplus}}\;\mathfrak{g}_{j}\,,
\end{equation*}
where each \,$\mathfrak{g}_{j} \subset \mathfrak{g}$\, is a simple Lie algebra.
\end{theorem}

          %%%%% ~~~~~~~~~~~~~~~~~~~~ %%%%%

\vskip 0.5cm
\begin{proposition}[Reductive $\Longrightarrow$ ``unitary'' inner product exists on compact real form]
\label{ReductiveImpliesUnitaryInnerProductOnRealForm}
\mbox{}
\vskip 0.1cm
\noindent
Suppose \,$\mathfrak{g} \,=\, \mathfrak{k}_{\C}$\, is a reductive complex Lie algebra,
with compact real form \,$\mathfrak{k} \,\subset\, \mathfrak{g}$.\,
Then, the following statements are true:
\begin{enumerate}
\item
    There exists an $\Re$-valued inner product
    \,$\langle\,\cdot\,,\,\cdot\,\rangle : \mathfrak{k} \times \mathfrak{k} \longrightarrow \Re$\,
    on \,$\mathfrak{k}$\,
    with respect to which the adjoint action of
    \,$\mathfrak{k}$\, on \,$\mathfrak{g}$\,
    is ``unitary'' in the sense that
    \begin{equation}\label{almostUnitarity}
    \langle\,\ad_{X}(Y)\,,\,Z\,\rangle
    \;\; = \;\;
        -\,\langle\,Y\,,\,\ad_{X}(Z)\,\rangle\,,
    \quad
    \textnormal{for each \,$X \in \mathfrak{k}$\, and each \,$Y,Z \in \mathfrak{g}$}
    \end{equation}
\item
    If we define an operation
    \,$\mathfrak{g} \longrightarrow \mathfrak{g} : X \longmapsto X^{*}$\,
    by
    \begin{equation*}
    \left(\,\overset{{\color{white}.}}{X_{1} + \i\cdot X_{2}}\,\right)^{*}
    \;\; := \;\;
        -\,X_{1} + \i\cdot X_{2}\,,
    \quad
    \textnormal{for each \,$X_{1}, X_{2} \in \mathfrak{k}$}\,,
    \end{equation*}
    then any inner product on \,$\mathfrak{g}$\, satisfying \eqref{almostUnitarity}
    also satisfies:
    \begin{equation}
    \langle\,\ad_{X}(Y)\,,\,Z\,\rangle
    \;\; = \;\;
        \langle\,Y\,,\,\ad_{X^{*}}(Z)\,\rangle\,,
    \quad
    \textnormal{for each \,$X,Y,Z \in \mathfrak{g}$}
    \end{equation}
\end{enumerate}
\end{proposition}
\proof
\begin{enumerate}
\item
    By hypothesis, \,$\mathfrak{k}$\, is the Lie algebra of a compact (real) matrix Lie group \,$K$.\,
    Next, recall that every complex representation of a compact Lie group is
    unitary\footnote{If $G$ is a compact Lie group and $\pi : G \longrightarrow \GL(V)$
    is a finite-dimensional complex representation,
    then there exists a $G$-invariant Hermitian inner product
    $\langle\,\cdot\,,\,\cdot\,\rangle$ on $V$,
    i.e., $\langle\,\pi(g)\cdot v \,,\,\pi(g)\cdot w\,\rangle \,=\, \langle\,v\,,\,w\,\rangle$,\,
    for each $g \in G$, and each $v, w \in V$.}; see, for example, Lemma 5.6, p.100, \cite{Sepanski2010}.
    Hence, there exists an ($\Re$-valued) inner product
    \,$\langle\,\cdot\,,\,\cdot\,\rangle : \mathfrak{k} \times \mathfrak{k} \longrightarrow \Re$\,
    that is (a) invariant with respect to the adjoint action
    \,$\Ad : K \longrightarrow \GL(\mathfrak{k})$,\, and (b) skew-symmetric with respect to
    the infinitesimal adjoint
    action\footnote{(differential of $\Ad : K \longrightarrow \GL(\mathfrak{k})$)}
    \,$\ad : \mathfrak{k} \longrightarrow \gl(\mathfrak{k})$,\, i.e.,
    \begin{equation}
    \langle\,\Ad(k) \cdot Y\,,\,\Ad(k) \cdot Z\,\rangle 
    \;\; = \;\;
        \langle\,Y\,,\,Z\,\rangle\,,
    \quad
    \textnormal{for each \,$k \in K$\, and each \,$Y,Z \in \mathfrak{k}$}
    \end{equation}
    and
    \begin{equation}
    \langle\,\ad_{X}(Y)\,,\,Z\,\rangle \,+\, \langle\,Y\,,\,\ad_{X}(Z)\,\rangle
    \;\; = \;\;
        0\,,
    \quad
    \textnormal{for each \,$X,Y,Z \in \mathfrak{k}$}
    \end{equation}
    Extension of \,$\ad_{X}$\, by {\color{red}$\C$-linearity} and
    extension of \,$\langle\,\cdot\,,\,\cdot\,\rangle$\, by {\color{red}bi-additivity}
    together yield the desired equality.
    Indeed, for each \,$X \in \mathfrak{k}$\, and each \,$Y, Z \in \mathfrak{g}$,\,
    where \,$Y = Y_{1} + \i\,Y_{2}$\, and \,$Z = Z_{1} + \i\,Z_{2}$\,
    with \,$Y_{1}, Y_{2}, Z_{1}, Z_{2} \in \mathfrak{k}$,\, we have:
    \begin{eqnarray*}
    \langle\,\ad_{X}(Y)\,,\,Z\,\rangle
    & = &
        \langle\,\ad_{X}(Y_{1}+\i\,Y_{2})\,,\,Z_{1} + \i\,Z_{2}\,\rangle
    \;\; = \;\;
        \langle\,\ad_{X}(Y_{1})+\ad_{X}(\i\,Y_{2})\,,\,Z_{1} + \i\,Z_{2}\,\rangle
    \\
    & = &
        \langle\,\ad_{X}(Y_{1})\,,\,Z_{1}\,\rangle
        \,+\,
        \langle\,\ad_{X}(Y_{1})\,,\,\i\,Z_{2}\,\rangle
        \,+\,
        \langle\,\ad_{X}(\i\,Y_{2})\,,\,Z_{1}\,\rangle
        \,+\,
        \langle\,\ad_{X}(\i\,Y_{2})\,,\,\i\,Z_{2}\,\rangle
    \\
    & = &
        -\,\langle\,Y_{1}\,,\,\ad_{X}(Z_{1})\,\rangle
        \,-\,\langle\,Y_{1}\,,\,\ad_{X}(\i\,Z_{2})\,\rangle
        \,-\,\langle\,\i\,Y_{2}\,,\,\ad_{X}(Z_{1})\,\rangle
        \,-\,\langle\,\i\,Y_{2}\,,\,\ad_{X}(\i\,Z_{2})\,\rangle
    % \\
    % & = &
    %     -\,\langle\,Y_{1}\,,\,\ad_{X}(Z_{1})\,\rangle
    %     -\,\langle\,Y_{1}\,,\,\ad_{X}(\i\,Z_{2})\,\rangle
    %     -\,\langle\,\i\,Y_{2}\,,\,\ad_{X}(Z_{1})\,\rangle
    %     -\,\langle\,\i,Y_{2}\,,\,\ad_{X}(\i\,Z_{2})\,\rangle
    \\
    & = &
        -\,\langle\,Y_{1}\,,\,\ad_{X}(Z_{1}+\i\,Z_{2})\,\rangle
        -\,\langle\,\i\,Y_{2}\,,\,\ad_{X}(Z_{1}+\i\,Z_{2})\,\rangle
    \\
    & = &
        -\,\langle\,Y_{1}+\i\,Y_{2}\,,\,\ad_{X}(Z_{1}+\i\,Z_{2})\,\rangle
    \\
    & = &
        -\,\langle\,Y\,,\,\ad_{X}(Z)\,\rangle\,,
    \quad
    \textnormal{as desired}
    \end{eqnarray*}
\item
    Extension of \,$\ad_{X}$\, by {\color{red}$\C$-linearity} and
    extension of
    \,$\langle\,\cdot\,,\,\cdot\,\rangle$\, by {\color{red}$\C$-bilinearity}
    (anti-linear in the first argument)
    together yield:
    \begin{eqnarray*}
    \langle\,\ad_{X}(Y)\,,\,Z\,\rangle
    & = &
        \langle\,\ad_{X_{1} + \i\,X_{2}}(Y)\,,\,Z\,\rangle
    \;\; = \;\;
        \langle\,\ad_{X_{1}}(Y) + \i\,\ad_{X_{2}}(Y)\,,\,Z\,\rangle
    \\
    & = &
        \langle\,\ad_{X_{1}}(Y)\,,\,Z\,\rangle
        \; + \;
        ({\color{red}-\,\i})\,\langle\,\ad_{X_{2}}(Y)\,,\,Z\,\rangle
    \;\; = \;\;
        {\color{red}-}\,\langle\,Y\,,\,\ad_{X_{1}}(Z)\,\rangle
        \; {\color{red}+} \;
        {\color{red}\i}\,\langle\,Y\,,\,\ad_{X_{2}}(Z)\,\rangle
    \\
    & = &
        \langle\,Y\,,\,\ad_{({\color{red}-}X_{1})}(Z)\,\rangle
        \; + \;
        \langle\,Y\,,\,\ad_{({\color{red}\i}X_{2})}(Z)\,\rangle
    \;\; = \;\;
        \langle\,Y\,,\,\ad_{(-X_{1}+\i X_{2})}(Z)\,\rangle
    \\
    & = &
        \langle\,Y\,,\,\ad_{X^{*}}(Z)\,\rangle\,,
    \end{eqnarray*}
    as required.
    \qed
\end{enumerate}


          %%%%% ~~~~~~~~~~~~~~~~~~~~ %%%%%

          %%%%% ~~~~~~~~~~~~~~~~~~~~ %%%%%

          %%%%% ~~~~~~~~~~~~~~~~~~~~ %%%%%

          %%%%% ~~~~~~~~~~~~~~~~~~~~ %%%%%

          %%%%% ~~~~~~~~~~~~~~~~~~~~ %%%%%

          %%%%% ~~~~~~~~~~~~~~~~~~~~ %%%%%

          %%%%% ~~~~~~~~~~~~~~~~~~~~ %%%%%
